\chapter{Overdetermination}
\label{ch:overdetermination}

Level is not the only way a soundtrack claims attention. An abrupt onset claims it; so does a rhythmic low-frequency pulse, a rising swell, and bright, noisy spectral content. This chapter considers what happens when several of these channels are used at once, and what the present measurements can and cannot say about it.

\section{The idea}

A single moment in an action sequence may contain a physical impact, a vocal cry, a percussive hit in the score, a rising orchestral or synthetic swell, and the broadband noise of breaking glass or gunfire. Each layer, taken alone, is ordinary craft. Murch's account of density describes how layers are distributed so that each remains legible \citep{murch2005}, and Chion's notion of added value describes how sound directs the interpretation of the image \citep{chion1994}. The concern here is different. When several channels are elevated together, they may be communicating the same thing, an instruction to attend now, through redundant means. The word \emph{overdetermination} is used for this: an effect produced by more causes than are needed to produce it.

Redundancy is not in itself a defect; it makes a signal robust. The hypothesis of Chapter~\ref{ch:framework} is narrower. It is that, at a given level of exceedance, a moment in which more salience channels are elevated is more likely to provoke a volume reduction, and that sustained overdetermination contributes to the listener's sense that the soundtrack never stops demanding attention.

\section{Measurement}

The overdetermination index $\Omega(t)$ counts, for each second, how many of five channels are elevated: level exceedance ($E\ge\delta$), transient barrage, low-frequency onset activity, harshness, and an active swell event. Channels other than exceedance count as elevated when their per-second score exceeds 1.5 on the pipeline's standardised scale.

Because several channels share components with level (see below), the index is also computed in a residualised form, $\Omega_{\text{res}}$. For each channel, a quadratic in $L(t)$ is fitted across audible seconds, and the channel counts as elevated only if it exceeds the raw threshold \emph{and} at least half that margin remains after the level-predicted value is subtracted. By construction $\Omega_{\text{res}}\le\Omega$.

\begin{table}[htbp]\centering\small
\begin{tabular}{lrr}
\toprule
 & $\Omega$ & $\Omega_{\text{res}}$ \\
\midrule
Mean within occupation blocks & \resOmegaIn & \resOmegaResIn \\
Mean outside occupation blocks & \resOmegaOut & \resOmegaResOut \\
Share of audible seconds with $\ge 3$ channels & \resOmegaThreeShare & \resOmegaResThreeShare \\
Share of block seconds with $\ge 3$ channels & \resOmegaThreeShareIn & \resOmegaResThreeShareIn \\
\bottomrule
\end{tabular}
\caption{Overdetermination, raw and residualised on level.}\label{tab:omega}
\end{table}

The difference between the two columns estimates how much of the apparent co-elevation is mechanical. What remains in the residualised column is co-elevation that level does not explain.

\section{Result for the case}

At first sight the index doubles inside occupation blocks, raw and residualised alike. But one of its five channels is exceedance itself, which is present in nearly every block second by definition. Removing that term leaves the number of \emph{other} channels elevated per second: \xOtherRawIn{} inside blocks against \xOtherRawOut{} outside on the raw scale, and \xOtherResIn{} against \xOtherResOut{} after residualising on level. The doubling is carried entirely by exceedance. Excluding it, the blocks are no more overdetermined than the rest of the film, and after removal of the level dependence slightly less so. Seconds with three or more channels elevated are rare everywhere: \resOmegaResThreeShare{} of audible seconds and \resOmegaResThreeShareIn{} of block seconds on the residualised index.

The channel means tell the same story with one qualification. Low-frequency onset activity is much higher on average inside blocks (\xDrumIn{} against \xDrumOut), and barrage somewhat higher (\xBarrageIn{} against \xBarrageOut), while harshness is about the same (\xHarshIn{} against \xHarshOut). These elevations are mostly below the channel threshold. The blocks carry more low-frequency rhythmic activity, but not in a way that stacks with the other channels.

For this film, at these thresholds, the soundtrack-side premise of the overdetermination hypothesis is not supported. What distinguishes its occupied stretches is sustained level, with somewhat denser low-frequency activity, not a pile-up of independent salience channels.

\section{What the index cannot show yet}

Three limitations prevent these numbers from testing the hypothesis.

First, the channels are not independent by construction. The barrage and harshness scores each include a level term, and every onset-based measure is audibility-weighted by level. A rise in level therefore raises several channels mechanically. The residualised index addresses the part of this coupling that a smooth function of level can capture.\footnote{Residualising on a quadratic in level removes the average dependence of each channel on level, not every interaction. A channel whose relation to level differs between scenes (for example, harshness that rises with level in gunfire but not in music) is only partly corrected.}

Second, the channels are acoustic, not semantic. Elevated low-frequency onset activity may be a drum, an engine, or a sequence of impacts. The overdetermination hypothesis is about distinct sources carrying the same message, which requires knowing what the sources are. The audition clips exist for this purpose, and Chapter~\ref{ch:response} describes how their annotation should be conducted.

Third, the hypothesis is about behaviour. Co-elevation is only interesting if it predicts volume reductions beyond what exceedance predicts, and that requires the response data specified in the next chapter.

\section{A note on the earlier reading}

Version~3's plots appeared to show every channel rising together between about 81 and 84 minutes, and an early reading took this as the clearest instance of overdetermination in the film. Chapter~\ref{ch:failures} showed that this stretch was quiet and that its apparent intensity was an artefact of level-invariant features. The episode is a warning specific to this chapter's subject: an index that counts coincident elevations will find the most coincidences wherever a shared artefact touches every channel at once. Audibility weighting removes the artefact that produced that particular coincidence. Residualising on level removes much of the mechanical coupling that remains.

\section{Status}

Overdetermination is graded as hypothesised. For this film the index finds no excess co-elevation in occupation blocks once exceedance itself is set aside. That result is limited by the acoustic channels chosen and their thresholds, and it says nothing yet about sources or behaviour: a scream, a gunshot, and a drum hit in the same second could all register in one channel. The hypothesis is therefore weakened rather than rejected. Source labels from the blinded annotation are the next test of it.
